\section{Peripheral Subsystems}
\label{sec:peripherals}

NavHAL provides a comprehensive suite of peripheral drivers, each engineered with a focus on ease-of-use, performance, and hardware independence. By leveraging modern C features and efficient memory-mapped I/O techniques, NavHAL offers a premium developer experience comparable to high-level frameworks but with the efficiency of bare-metal code.

\subsection{GPIO: The Universal Pin Interface}
The GPIO subsystem is the foundation of NavHAL's interaction with the physical world. It abstracts pins into a unified format (e.g., \texttt{GPIO\_PA05}) while providing granular control over electrical characteristics.

\subsubsection{Key Features}
\begin{itemize}
    \item \textbf{Mode Management}: Simple API for switching between Input, Output, Alternate Function, and Analog modes.
    \item \textbf{Electrical Configuration}: Support for internal Pull-Up/Pull-Down (PUPD) resistors, Output Type (Push-Pull vs. Open-Drain), and configurable Slew Rates.
    \item \textbf{Type-Safe Pin Mapping}: Pins are represented as unique identifiers, preventing port-mismatch errors during configuration.
\end{itemize}

\begin{lstlisting}[language=C, caption=GPIO Initialization and Control]
hal_gpio_setmode(GPIO_PA05, GPIO_OUTPUT, GPIO_PUPD_NONE);
hal_gpio_digitalwrite(GPIO_PA05, GPIO_HIGH);
\end{lstlisting}

\subsection{UART: High-Performance Serial Communication}
The UART module in NavHAL is a standout feature, offering advanced abstractions that simplify complex serial protocols.

\subsubsection{Type-Generic Write Macros}
One of NavHAL's most powerful features is the use of the C11 \texttt{\_Generic} keyword to provide a unified \texttt{write} interface. The \texttt{uart\_write()} macros automatically select the correct internal function based on the data type passed (char, int, float, or string).

\begin{lstlisting}[language=C, caption=Type-Generic UART Transmission]
uart2_write("Temperature: "); 
uart2_write(25.5f);           // Automatically calls uart2_write_float
uart2_write('\n');            // Automatically calls uart2_write_char
\end{lstlisting}

\subsubsection{DMA Integration}
For data-intensive applications, NavHAL provides a seamless path to hardware-accelerated communication. The UART driver includes built-in support for Direct Memory Access (DMA), allowing background transmissions and circular-buffer reception that offloads the CPU entirely.

\subsection{I2C and SPI: Synchronous Bus Abstractions}
The synchronous communication modules (I2C and SPI) follow a consistent transaction-based model.

\begin{itemize}
    \item \textbf{I2C Controller}: Supports Master/Slave modes, multi-byte transactions, and integrated DMA for high-bandwidth sensor data retrieval.
    \item \textbf{SPI Interface}: High-speed full-duplex communication with configurable CPOL/CPHA settings, ideal for driving displays and external storage modules.
\end{itemize}

\subsection{Timers and Pulse Width Modulation (PWM)}
NavHAL abstracts the complex timer hardware into two primary services:
\begin{enumerate}
    \item \textbf{Timebase Services}: Provides precise delays and periodic interrupts via the Systick or hardware timers.
    \item \textbf{PWM Engine}: Simplifies signal generation for motor control, LED dimming, and audio applications with an intuitive \texttt{hal\_pwm\_set\_duty()} API.
\end{enumerate}

\subsection{DMA Controller: The Performance Backbone}
The Direct Memory Access (DMA) subsystem is a critical component of NavHAL's high-performance architecture. It acts as a transparent data mover between peripherals and memory, ensuring that data-heavy tasks (like UART logging or sensor polling) do not stall the main execution thread. NavHAL's DMA management includes:
\begin{itemize}
    \item \textbf{Channel Mapping}: Automatic resolution of DMA streams and channels for supported peripherals.
    \item \textbf{Interrupt Integration}: Optional callback support for completion and error events.
\end{itemize}
